% Interaction gravitationnelle Terre-Lune : deux forces opposées (corrigé de l'exercice 3, chapitre 11)
\begin{tikzpicture}[x=1cm,y=1cm,force/.style={otGreenDark,line width=1.2pt,-{Stealth[length=2.4mm]}}]
\shade[ball color=otNavy!35] (0,0) circle (0.9); \node[pcLegendeG] at (0,0) {Terre};
\shade[ball color=otMuted!40] (8,0) circle (0.3); \node[pcLegende,anchor=north] at (8,-0.35) {Lune};
\draw[force] (0.9,0) -- (2.4,0) node[pcLegende,anchor=south,text=otGreenDark] {$\vec{F}_{\mathrm{L}/\mathrm{T}}$};
\draw[force] (7.7,0) -- (6.2,0) node[pcLegende,anchor=south,text=otGreenDark] {$\vec{F}_{\mathrm{T}/\mathrm{L}}$};
\draw[otMuted,line width=0.4pt,dash pattern=on 2pt off 1.5pt] (0,0) -- (0.9,0) (7.7,0) -- (8,0);
\draw[otMuted,line width=0.4pt,{Stealth[length=1.4mm]}-{Stealth[length=1.4mm]}] (0,-1.2) -- (8,-1.2) node[midway,below,pcLegende] {$d_\mathrm{TL}=\qty{3,84e8}{m}$ (échelle non respectée)};
\end{tikzpicture}
\par\smallskip\legende{Les deux forces ont la même valeur, \qty{1,98e20}{N}, et des sens opposés : la Lune attire la Terre autant que la Terre attire la Lune.}
